\chapter*{How the book is organised}
\addcontentsline{toc}{chapter}{How the book is organised}

This book is a guided journey through the ideas needed to understand, design, build, and evaluate digital systems. It is not a timetable or a list of institutional requirements. The chapters are arranged so that each part supplies concepts that make the next part easier, while still allowing readers to move directly to a topic that matters to them.

\begingroup
\setlength{\tabcolsep}{4pt}
\begin{longtable}{p{0.20\textwidth}p{0.12\textwidth}p{0.29\textwidth}p{0.29\textwidth}}
\caption{The architecture of the book.}\label{tab:book-map}\\
\toprule
\textbf{Part or focus} & \textbf{Chapters} & \textbf{Central questions} & \textbf{What the reader develops}\\
\midrule
\endfirsthead
\toprule
\textbf{Part or focus} & \textbf{Chapters} & \textbf{Central questions} & \textbf{What the reader develops}\\
\midrule
\endhead
Foundations & 1--4 & What is a digital system? How can a problem be represented computationally? & Problem statements, models, programs, algorithms, and computational reasoning\\
Engineering systems & 5--10 & How can software be designed, implemented, stored, tested, secured, and operated? & Requirements, architecture, code, data, testing, security, and operational judgment\\
People and practice & 11--14 & How do people experience systems, and how do organisations and teams shape them? & Interaction design, evaluation, organisational analysis, agile practice, and collaboration\\
Evidence and independent work & 15--18 & How can claims be investigated, evaluated, and connected to responsible action? & Research design, human-centred AI, experimentation, entrepreneurship, and thesis reasoning\\
Integration projects & 19--20 & What happens when the concepts meet in a complete problem? & A worked case, project ideas, evaluation plans, and final synthesis\\
Appendices and tools & A--D & Which mathematical, technical, and conceptual tools support the main chapters? & Mathematics, Python and Git, reference tables, and a glossary\\
\bottomrule
\end{longtable}
\endgroup

The parts are not prerequisites in the bureaucratic sense, and they are not meant to be completed at a fixed speed. A beginner may read sequentially and run the examples as they appear. An experienced reader may begin with the case study, research chapter, security material, or a specific engineering problem, then return to earlier chapters when a foundation is needed.

\section*{The recurring movement}
The book repeatedly moves through a chain of questions:
\[
\text{problem} \rightarrow \text{model} \rightarrow \text{program} \rightarrow \text{component} \rightarrow \text{system} \rightarrow \text{organisation} \rightarrow \text{evidence}.
\]
This is not a one-way lifecycle. User research can change the problem; production incidents can change the model; organisational constraints can change the architecture; and evidence can show that an intervention did not achieve its intended effect. The purpose of the sequence is therefore not to prescribe a recipe, but to make the connections visible.
